%%--------------------------------------------------------------------
%1--------------------------------------------------------------------
\vspace*{\fill}
\begin{glyph}{Grove (林)}{grove}\label{glyph:grove}
Grove.
\end{glyph}

\begin{comps}
\comprtwocone &  木 + 木 \\\hline
    \comprthreecone & Tree (\ref{glyph:tree}) + Tree (\ref{glyph:tree}) \\\hline
\end{comps}

\begin{remglyph}
--- \\\hline
\end{remglyph}

\begin{prons}
\pronsrtwocone & \textbf{リン (RIN)} & \textbf{はやし (hayashi)}\\\hline
\pronsrthreecone & --- & --- \\\hline
\end{prons}

\begin{remprons}
    The \hooglig{grove} is \comon{leaning} towards \comkun{high-ash-investments}.\\\hline
\end{remprons}

%{\middel}
% &  ()\\\hline
\begin{somme}
    \rye{3}{As the common Japanese family name in the second example below
    shows, this reading becomes voiced when not in the first position.}
    林  & \textit{grove}  & はやし (hayashi)  \\\hline
    小林さん  & small + grove = \textit{Ko{\middel}bayashi-san}  & こ{\middel}ばやし{\middel}さん
    (ko{\middel}bayashi{\middel}san)  \\\hline
    林道  & grove + road = \textit{woodland trail; road}  & リン{\middel}ドウ (RIN{\middel}D\={O})  \\\hline
    人工林  & person + craft + grove = \textit{planted forest}  & ジン{\middel}コウ{\middel}リン
    (JIN{\middel}K\={O}{\middel}RIN)  \\\hline
    公有林  & public + have + grove = \textit{public woodland}  & コウ{\middel}ユウ{\middel}リン
    (K\={O}{\middel}Y\={U}{\middel}RIN)  \\\hline
\end{somme}

\begin{sentence}
    \underline{あの}{\;}\underline{\ruby{国}{くに}}に\ruby{は}{わ}\underline{\ruby{公有林}{コウユウリン}}が\underline{\ruby{少}{すく}ない}。\\\hline
    ano kuni ni wa K\={O}{\middel}Y\={U}{\middel}RIN ga suku{\middel}nai.\\\hline
    (that) (country) (public woodland) (few)\\\hline
    =\textit{That country has little public woodland.}\\\hline
\end{sentence}
\end{multicols}
\vspace*{-\baselineskip}
\vhrulefill{5pt}
%%--------------------------------------------------------------------
%2--------------------------------------------------------------------
\begin{multicols}{2}
\begin{glyph}{Forest (森)}{forest}\label{glyph:forest}
Forest.
\end{glyph}

\begin{comps}
\comprtwocone & 木 + 木 + 木 \\\hline
    \comprthreecone & Tree (\ref{glyph:tree}) + Tree (\ref{glyph:tree}) + Tree
    (\ref{glyph:tree}) \\\hline
\end{comps}

\begin{remglyph}
---\\\hline
\end{remglyph}

\begin{prons}
\pronsrtwocone & \textbf{シン (SHIN)} & \textbf{もり (mori)}\\\hline
\pronsrthreecone & --- & --- \\\hline
\end{prons}

\begin{remprons}
    A \hooglig{forest} with a \comkun{moribund} \comon{sheen}.\\\hline
\end{remprons}

%{\middel}
% &  ()\\\hline
\begin{somme}
    森  & \textit{forest}  & もり (mori)  \\\hline
    森林  & forest + grove = \textit{forests}  & シン{\middel}リン (SHIN{\middel}RIN)  \\\hline
    森林学  & forest + grove + study = \textit{forestry}  & シン{\middel}リン{\middel}ガク
    (SHIN{\middel}RIN{\middel}GAKU)  \\\hline
\end{somme}

\begin{sentence}
    \underline{\ruby{森}{もり}}に\underline{\ruby{入}{はい}る}{\;}\underline{\ruby{時}{とき}}\ruby{は}{わ}\underline{\ruby{注意}{チュウイ}}{\;}\underline{して}{\;}\underline{\ruby{下}{くだ}さい}。\\\hline
    mori ni hai{\middel}ru toki wa CH\={U}{\middel}I shite kuda{\middel}sai.\\\hline
    (forest) (enter) (time) (attention) (do please)\\\hline
    =\textit{Please be careful when going into the forest.}\\\hline
\end{sentence}
\end{multicols}
\vspace*{-\baselineskip}
\vhrulefill{5pt}
\vspace*{\fill}
\pagebreak
%%--------------------------------------------------------------------
%3--------------------------------------------------------------------
\vspace*{\fill}
\begin{multicols}{2}
\begin{glyph}{Dawn (旦)}{dawn}\label{glyph:dawn}
Dawn.\\\hline
This kanji features more often as a component in other kanji.
\end{glyph}

\begin{comps}
\comprtwocone & 日 + 一 \\\hline
    \comprthreecone & Sun (\ref{glyph:sun}) + One (\ref{glyph:one}) \\\hline
\end{comps}

\begin{remglyph}
    \hooglig{Dawn} is stage \remcom{one} of the \remcom{day}.\\\hline
\end{remglyph}

\begin{prons}
\pronsrtwocone & \textbf{タン (TAN)} & --- \\\hline
    \pronsrthreecone & ダン (DAN) & --- \\\hline
\end{prons}

\begin{remprons}
    \hooglig{Dawn} \ucomon{dunking} by the \comon{tons}.\\\hline
\end{remprons}

%{\middel}
% &  ()\\\hline
\begin{somme}
    一旦  & once + dawn = \textit{once}  & イッ{\middel}タン (IT{\middel}TAN)  \\\hline
    元旦  & basis + dawn = \textit{New Year's Day}  & ガン{\middel}タン (GAN{\middel}TAN)  \\\hline
\end{somme}

\begin{sentence}
    \underline{\ruby{元旦}{ガンタン}}に\underline{お\ruby{寺}{てら}}\ruby{へ}{え}\underline{\ruby{行}{い}きました}。\\\hline
GAN{\middel}TAN ni o{\middel}tera e i{\middel}kimashita.\\\hline
    (New Year's Day) (temple) (went)\\\hline
    =\textit{(I) went to a temple on New Year's Day.}\\\hline
\end{sentence}
\end{multicols}
\vspace*{-\baselineskip}
\vhrulefill{5pt}
%%--------------------------------------------------------------------
%4--------------------------------------------------------------------
\begin{multicols}{2}
\begin{glyph}{Black (黒)}{black}\label{glyph:black}
Black.
\end{glyph}

\begin{comps}
\comprtwocone & 里 + 灬 \\\hline
    \comprthreecone & Black (\ref{glyph:black}) + Fire (\ref{glyph:fire}) \\\hline
\end{comps}

\begin{remglyph}
    \hooglig{Black} \remcom{fire} consumed the \remcom{hamlet}.\\\hline
\end{remglyph}

\begin{prons}
\pronsrtwocone & \textbf{コク (KOKU)} & \textbf{くろ (kuro)}\\\hline
\pronsrthreecone & --- & --- \\\hline
\end{prons}

\begin{remprons}
    \hooglig{Black} \comkun{cool objects} \comon{cocked} my interests.\\\hline
\end{remprons}

%{\middel}
% &  ()\\\hline
\begin{somme}
    \rye{3}{With the exception of the third example, the kun-yomi is always
    voiced in the second position (as it appears in the fourth compound).}
    黒  & \textit{black (noun)}  & くろ (kuro)  \\\hline
    黒い  & \textit{black (adjective)}  & くろ{\middel}い (kuro{\middel}i)  \\\hline
    白黒  & white + black = \textit{black-and-white}  & しら{\middel}くろ (shira{\middel}kuro)  \\\hline
    赤黒い  & red + black = \textit{dark red}  & あか　ぐろ{\middel}い (aka guro{\middel}i)  \\\hline
    黒人  & black + person = \textit{a black person}  & コク{\middel}ジン (KOKU{\middel}JIN)  \\\hline
\end{somme}

\begin{sentence}
    \underline{この}{\;}\underline{\ruby{黒}{くろ}い}{\;}\underline{\ruby{馬}{うま}}\ruby{は}{わ}\underline{\ruby{美}{うつく}しい}{\;}\underline{ですね}。\\\hline
kono kuro{\middel}i uma wa utsuku{\middel}shii desu ne.\\\hline
    (this) (black) (horse) (beautiful) (isn't it)\\\hline
    =\textit{This black horse is beautiful, isn't it?}\\\hline
\end{sentence}
\end{multicols}
\vspace*{-\baselineskip}
\vhrulefill{5pt}
\vspace*{\fill}
\pagebreak
%%--------------------------------------------------------------------
%5--------------------------------------------------------------------
\vspace*{\fill}
\begin{multicols}{2}
\begin{glyph}{Mix (交)}{mix}\label{glyph:mix}
A general sense of things mixing in both physical and figurative ways.
    ``Associating'' with people or ``exchanging'' words are examples of the
    latter.
\end{glyph}

\begin{comps}
\comprtwocone & 亠 + 父 \\\hline
    \comprthreecone & head/above/lid + Father (\ref{glyph:father}) \\\hline
\end{comps}

\begin{remglyph}
    \hooglig{Mix} up \remcom{father's} \remcom{head}.\\\hline
\end{remglyph}

\begin{prons}
    \pronsrtwocone & \textbf{コウ (K\={O})} & \textbf{ま (ma)}\\\hline
    \pronsrthreecone & --- & まじ、か (maji, ka) \\\hline
\end{prons}

\begin{remprons}
    A \hooglig{mixup} \comon{caused} by a \ucomkun{customer} \comkun{muddying}
    \ucomkun{much information}.\\\hline
\end{remprons}

%{\middel}
% &  ()\\\hline
\begin{somme}
    交ざる (intr.)  & \textit{to mingle}  & ま{\middel}ざる (ma{\middel}zaru)  \\\hline
    交ぜる (tr.)  & \textit{to mix (something)}  & ま{\middel}ぜる (ma{\middel}zeru)  \\\hline
    外交  & outside + mix = \textit{foreign policy}  & ガイ{\middel}コウ (GAI{\middel}K\={O})  \\\hline
\end{somme}

\begin{sentence}
    \underline{あの}{\;}\underline{\ruby{国}{くに}}の\underline{\ruby{外交}{ガイコウ}}が\underline{\ruby{分}{わ}かりません}。\\\hline
    ano kuni no GAI{\middel}K\={O} ga wa{\middel}karimasen.\\\hline
    (that) (country) (foreign policy) (don't understand)\\\hline
    =\textit{(I) don't understand that country's foreign policy.}\\\hline
\end{sentence}
\end{multicols}
\vspace*{-\baselineskip}
\vhrulefill{5pt}
%%--------------------------------------------------------------------
%6--------------------------------------------------------------------
\begin{multicols}{2}
\begin{glyph}{Not (不)}{not}\label{glyph:not}
Expressing the idea of ``not'' or ``un-''.  This is an important negating prefix
in Japanese.
\end{glyph}

\begin{comps}
\comprtwocone & 一 + 亻 + 丶 \\\hline
    \comprthreecone & One (\ref{glyph:one}) + Person (\ref{glyph:person}) + dot \\\hline
\end{comps}

\begin{remglyph}
    \hooglig{Not} \remcom{one} \remcom{person} without a \remcom{dot} of sin.\\\hline
\end{remglyph}

\begin{prons}
\pronsrtwocone & \textbf{フ (FU)} & --- \\\hline
    \pronsrthreecone & ブ (BU) & --- \\\hline
\end{prons}

\begin{remprons}
    \hooglig{Not} a \comon{foofy} \ucomon{bully}.\\\hline
\end{remprons}

%{\middel}
% &  ()\\\hline
\begin{somme}
    \rye{3}{Note how this kanji can function as a prefix with complete words, as
    it does in the final example.}
    不意  & not + mind = \textit{unexpectedly}  & フ{\middel}イ (FU{\middel}I)  \\\hline
    不安  & not + ease = \textit{uneasy}  & フ{\middel}アン (FU{\middel}AN)  \\\hline
    不明  & not + bright = \textit{unclear}  & フ{\middel}メイ (FU{\middel}MEI)  \\\hline
    不死  & not + death = \textit{immortal}  & フ{\middel}シ (FU{\middel}SHI)  \\\hline
    不親切  & not + parent + cut = \textit{unkind}  & フ{\middel}シン{\middel}セツ (FU{\middel}SHIN{\middel}SETSU)  \\\hline
\end{somme}

\begin{sentence}
    \underline{\ruby{小林}{こばやし}さん}の\underline{\ruby{意図}{イト}}\ruby{は}{わ}\underline{\ruby{不明}{フメイ}}です。\\\hline
Ko{\middel}bayashi{\middel}san no I{\middel}TO wa FU{\middel}MEI desu.\\\hline
    (Kobayashi-san) (aim) (unclear)\\\hline
    =\textit{Kobayashi-san's aim is unclear.}\\\hline
\end{sentence}
\end{multicols}
\vspace*{-\baselineskip}
\vhrulefill{5pt}
\vspace*{\fill}
\pagebreak
%%--------------------------------------------------------------------
%7--------------------------------------------------------------------
\vspace*{\fill}
\begin{multicols}{2}
\begin{glyph}{After (後)}{after}\label{glyph:after}
After/Later/Behind.  These meanings apply to the ideas of both physical location
and time.
\end{glyph}

\begin{comps}
\comprtwocone & 彳 + 幺 + 夊 \\\hline
\comprthreecone & walking slowly + small/young + go \\\hline
\end{comps}

\begin{remglyph}
\hooglig{Afterwards} I \remcom{go} with the \remcom{young} children
    \remcom{walking slowly}.\\\hline
\end{remglyph}

\begin{prons}
    \pronsrtwocone & \textbf{コウ、ゴ (K\={O}, GO)} & \textbf{あと、うし (ato,
    ushi)}\\\hline
    \rye{3}{Look for コウ in the first position; あと and うしろ (this reading
    is always accompanied by ろ), appear less frequently here.\\\hline
    ゴ pops up in only a few words.  In the second position, moreover, the
    reading will invariably be ゴ.}
    \pronsrthreecone & --- & のち、おく (nochi, oku) \\\hline
\end{prons}

\begin{remprons}
    \hooglig{After} \ucomkun{notching} your way up, \comkun{you ship} the
    \comon{gobs} to \comon{cause} the \comkun{atoll} to be \ucomkun{occupied}.\\\hline
\end{remprons}

%{\middel}
% &  ()\\\hline
\begin{somme}
    \rye{3}{It is worth having another look at the compounds for ``前'' in
    Chapter 7 (Entry 133), to see how the fourth, fifth, and sixth examples
    given there ``pair up'' with the three below.}
    後  & \textit{after}  & あと (ato)  \\\hline
    後ろ  & \textit{behind}  & うし{\middel}ろ (ushi{\middel}ro)  \\\hline
    後年  & after + year = \textit{in later years}  & コウ{\middel}ネン (K\={O}{\middel}NEN)  \\\hline
    後半  & after + half = \textit{second half}  & コウ{\middel}ハン (K\={O}{\middel}HAN)  \\\hline
    後者  & after + individual = \textit{the latter}  & コウ{\middel}シャ (K\={O}{\middel}SHA)  \\\hline
    午後  & noon + after = \textit{P.M.; afternoon}  & ゴ{\middel}ゴ (GO{\middel}GO)  \\\hline
最後  & most + after = \textit{the last}  & サイ{\middel}ゴ (SAI{\middel}GO)  \\\hline
\end{somme}

\begin{sentence}
    \underline{\ruby{後}{あと}}で\underline{\ruby{友人}{ユウジン}}{\;}\underline{と}{\;}\underline{\ruby{東京}{トウキョウ}}\ruby{へ}{え}\underline{\ruby{行}{い}きます}。\\\hline
    ato de Y\={U}{\middel}JIN to T\={O}{\middel}KY\={O} e i{\middel}kimasu.\\\hline
    (after) (friend) (with) (Tokyo) (go).\\\hline
    =\textit{Afterwards, I'll go to Tokyo with a friend.}\\\hline
\end{sentence}
\end{multicols}
\vspace*{-\baselineskip}
\vhrulefill{5pt}
%%--------------------------------------------------------------------
%8--------------------------------------------------------------------
\begin{multicols}{2}
\begin{glyph}{Most (最)}{most}\label{glyph:most}
Most.\\\hline
Note how the first stroke of ``耳'' has been lengthened.
\end{glyph}

\begin{comps}
\comprtwocone & 日 + 最 \\\hline
    \comprthreecone & Sun (\ref{glyph:sun}) + Take (\ref{glyph:take}) \\\hline
\end{comps}

\begin{remglyph}
    \hooglig{Mostly} the \remcom{sun} \remcom{takes} it all.\\\hline
\end{remglyph}

\begin{prons}
\pronsrtwocone & \textbf{サイ (SAI)} & --- \\\hline
    \pronsrthreecone & --- & もっと、も (motto, mo) \\\hline
\end{prons}

\begin{remprons}
    The \hooglig{most} \comon{psychedellic} \ucomkun{moral} \ucomkun{motto}.\\\hline
\end{remprons}

%{\middel}
% &  ()\\\hline
\begin{somme}
    最小  & most + small = \textit{smallest}  & サイ{\middel}ショウ (SAI{\middel}SH\={O})  \\\hline
    最少  & most + few = \textit{fewest}  & サイ{\middel}ショウ (SAI{\middel}SH\={O})  \\\hline
    最高  & most + tall = \textit{highest; the best}  & サイ{\middel}コウ (SAI{\middel}K\={O})  \\\hline
    最古  & most + old = \textit{oldest}  & サイ{\middel}コ (SAI{\middel}KO)  \\\hline
    最新  & most + new = \textit{newest}  & サイ{\middel}シン (SAI{\middel}SHIN)  \\\hline
    最後  & most + after = \textit{the last}  & サイ{\middel}ゴ (SAI{\middel}GO)  \\\hline
\end{somme}

\begin{sentence}
    \underline{この}{\;}\underline{\ruby{天気}{テンキ}}\ruby{は}{わ}\underline{\ruby{最高}{サイコウ}}{\;}\underline{です}よ。\\\hline
    kono TEN{\middel}KI wa SAI{\middel}K\={O} desu yo.\\\hline
    (this) (weather) (the best) (is)\\\hline
    =\textit{This weather is fantastic.}\\\hline
\end{sentence}
\end{multicols}
\vspace*{-\baselineskip}
\vhrulefill{5pt}
\vspace*{\fill}
\pagebreak
%%--------------------------------------------------------------------
%9--------------------------------------------------------------------
\vspace*{\fill}
\begin{multicols}{2}
\begin{glyph}{Fat (太)}{fat}\label{glyph:fat}
    Fat/Thick.  Note the difference between this character and ``犬'' (Entry
    171).
\end{glyph}

\begin{comps}
\comprtwocone & 大 + 丶 \\\hline
    \comprthreecone & Large (\ref{glyph:large}) + dot \\\hline
\end{comps}

\begin{remglyph}
    \hooglig{Fat} \remcom{dots} are \remcom{large}.\\\hline
\end{remglyph}

\begin{prons}
\pronsrtwocone & \textbf{タイ (TAI)} & \textbf{ふと (futo)}\\\hline
    \pronsrthreecone & タ (TA) & --- \\\hline
\end{prons}

\begin{remprons}
    The \hooglig{fat} \comkun{futon} was \comon{tightly} packed in the
    \ucomon{tub}.\\\hline
\end{remprons}

%{\middel}
% &  ()\\\hline
\begin{somme}
    太い  & \textit{fat; thick}  & ふと{\middel}い (futo{\middel}i)  \\\hline
    太る  & \textit{to grow fat}  & ふと{\middel}る (futo{\middel}ru)  \\\hline
    太子  & fat + child = \textit{crown prince}  & タイ{\middel}シ (TAI{\middel}SHI)  \\\hline
\end{somme}

\begin{sentence}
    \underline{\ruby{私}{わたし}}\ruby{は}{わ}\underline{\ruby{冬休}{ふゆやす}み}に\underline{とても}{\;}\underline{\ruby{太}{ふと}った}。\\\hline
watashi wa fuyu yasu{\middel}mi ni totemo futo{\middel}tta.\\\hline
    (I) (winter vacation) (really) (grew fat)\\\hline
    =\textit{I really put on weight during the winter vacation.}\\\hline
\end{sentence}
\end{multicols}
\vspace*{-\baselineskip}
\vhrulefill{5pt}
%%--------------------------------------------------------------------
%10-------------------------------------------------------------------
\begin{multicols}{2}
\begin{glyph}{Car (車)}{car}\label{glyph:car}
Car.  As the compounds below indicate, this character can also refer to a wide
    range of other vehicles.
\end{glyph}

\begin{comps}
\comprtwocone & 車 \\\hline
    \comprthreecone & Car/Cart (\ref{glyph:car}) \\\hline
\end{comps}

\begin{remglyph}
--- \\\hline
\end{remglyph}

\begin{prons}
\pronsrtwocone & \textbf{シャ (SHA)} & \textbf{くるま (kuruma)}\\\hline
\pronsrthreecone & --- & --- \\\hline
\end{prons}

\begin{remprons}
    The \hooglig{car} was \comon{shut} to hear a \comkun{cool rumour}.\\\hline
\end{remprons}

%{\middel}
% &  ()\\\hline
\begin{somme}
    車  & \textit{car; vehicle}  & くるま (kuruma)  \\\hline
    電車  & electric + car = \textit{(electric) train}  & デン{\middel}シャ (DEN{\middel}SHA)  \\\hline
    空車  & empty + car = \textit{(taxi) ``for hire''}  & クウ{\middel}シャ (K\={U}{\middel}SHA)  \\\hline
    車内  & car + inside = \textit{inside the car}  & シャ{\middel}ナイ (SHA{\middel}NAI)  \\\hline
    中古車  & middle + old + car = \textit{used car}  & チュウ{\middel}コ{\middel}シャ
    (CH\={U}{\middel}KO{\middel}SHA)  \\\hline
\end{somme}

\begin{sentence}
    \underline{\ruby{中村}{なかむら}さん}\ruby{は}{お}\underline{\ruby{中古車}{チュウコシャ}}\ruby{を}{お}\underline{\ruby{買}{か}いました}。\\\hline
    Naka{\middel}mura-san wa CH\={U}{\middel}KO{\middel}SHA o ka{\middel}imashita.\\\hline
    (Nakamura-san) (used car) (bought)\\\hline
    =\textit{Nakamura-san bought a used car.}\\\hline
\end{sentence}
\end{multicols}
\vspace*{-\baselineskip}
\vhrulefill{5pt}
\vspace*{\fill}
\pagebreak
%%--------------------------------------------------------------------
%11-------------------------------------------------------------------
\vspace*{\fill}
\begin{multicols}{2}
\begin{glyph}{Receive (受)}{receive}\label{glyph:receive}
Receive/Accept.  This interesting character encompasses a wide range of meanings
having to do with the idea of reception: ``taking'' an exam, ``undergoing'' an
    operation, or ``suffering'' injuries are several examples.
\end{glyph}

\begin{comps}
\comprtwocone & 爪 + 冖 + 又 \\\hline
    \comprthreecone & claw + cover/roof + Again (\ref{glyph:again}) \\\hline
\end{comps}

\begin{remglyph}
    \hooglig{Receive} the \remcom{clawed} \remcom{cover} \remcom{again}.\\\hline
\end{remglyph}

\begin{prons}
\pronsrtwocone & \textbf{ジュ (JU)} & \textbf{う (u)}\\\hline
\pronsrthreecone & --- & --- \\\hline
\end{prons}

\begin{remprons}
    \hooglig{Receive} \comon{judicious} \comkun{oomph}.\\\hline
\end{remprons}

%{\middel}
% &  ()\\\hline
\begin{somme}
\rye{3}{The third and fourth examples below show this to be another kanji
    sharing the pronounciation characteristics of 買.}
    受ける  & \textit{to receive; to  accept}  & う{\middel}ける (u{\middel}keru)  \\\hline
    受け取る  & receive + take = \textit{to receive}  & う{\middel}け　と{\middel}る (u{\middel}ke to{\middel}ru)  \\\hline
    受取  & receive + take = \textit{recipient}  & うけ{\middel}とり (uke{\middel}tori)  \\\hline
    受取人  & receive + take + person = \textit{recipient}  & うけ{\middel}とり{\middel}ニン
    (uke{\middel}tori{\middel}NIN)  \\\hline
    引き受ける  & pull + receive = \textit{to undertake}  & ひ{\middel}き　う{\middel}ける (hi{\middel}ki
    u{\middel}keru)  \\\hline
    受注  & recieve + pour = \textit{receipt of an order}  & ジュ{\middel}チュウ
    (JU{\middel}CH\={U})  \\\hline
\end{somme}

\begin{sentence}
    \underline{\ruby{田中}{たなか}さん}\ruby{は}{わ}\underline{\ruby{新}{あたら}しい}{\;}\underline{\ruby{本}{ホン}}\ruby{を}{お}\underline{\ruby{受}{う}け\ruby{取}{と}りました}。\\\hline
Ta{\middel}naka-san wa atara{\middel}shii HON o u{\middel}ke to{\middel}rimashita.\\\hline
    (Tanaka-san) (new) (book) (received)\\\hline
    =\textit{Tanaka-san received a new book.}\\\hline
\end{sentence}
\end{multicols}
\vspace*{-\baselineskip}
\vhrulefill{5pt}
%%--------------------------------------------------------------------
%12-------------------------------------------------------------------
\begin{multicols}{2}
\begin{glyph}{Evil (凶)}{evil}\label{glyph:evil}
Evil/misfortune.
\end{glyph}

\begin{comps}
\comprtwocone & 丿 + 丶 + 凵  \\\hline
\comprthreecone & slash + dot + wide open mouth \\\hline
\end{comps}

\begin{remglyph}
    \hooglig{Evil} is \remcom{dotted} with \remcom{slashes}, it leaves your
    \remcom{mouth wide open}.\\\hline
\end{remglyph}

\begin{prons}
    \pronsrtwocone & \textbf{キョウ (KY\={O})} & --- \\\hline
\pronsrthreecone & ---  & --- \\\hline
\end{prons}

\begin{remprons}
    \hooglig{Evil} \comon{Kyoto}.\\\hline
\end{remprons}

%{\middel}
% &  ()\\\hline
\begin{somme}
    凶年  & evil + year = \textit{a bad year}  & キョウ{\middel}ネン (KY\={O}{\middel}NEN)  \\\hline
    凶行  & evil + go = \textit{violence}  & キョウ{\middel}コウ (KY\={O}{\middel}K\={O})  \\\hline
\end{somme}

\begin{sentence}
    \underline{\ruby{凶年}{キョウネン}}だ\underline{そう}です。\\\hline
    KY\={O}{\middel}NEN da s\={o} desu.\\\hline
    (bad year) (seems)\\\hline
    =\textit{Apparently it's a bad year.}\\\hline
\end{sentence}
\end{multicols}
\vspace*{-\baselineskip}
\vhrulefill{5pt}
\vspace*{\fill}
\pagebreak
%%--------------------------------------------------------------------
%13-------------------------------------------------------------------
\vspace*{\fill}
\begin{multicols}{2}
\begin{glyph}{Sea (海)}{sea}\label{glyph:sea}
Sea.
\end{glyph}

\begin{comps}
\comprtwocone & 氵 + 毎 \\\hline
    \comprthreecone & Water [variant] (\ref{glyph:water}) + Every (\ref{glyph:every}) \\\hline
\end{comps}

\begin{remglyph}
    The \hooglig{sea} has \remcom{water} \remcom{everywhere}.\\\hline
\end{remglyph}

\begin{prons}
\pronsrtwocone & \textbf{カイ (KAI)} & \textbf{うみ (umi)}\\\hline
\pronsrthreecone & --- & --- \\\hline
\end{prons}

\begin{remprons}
    \hooglig{Sea} \comon{kite-surfing} is \comkun{booming}.\\\hline
\end{remprons}

%{\middel}
% &  ()\\\hline
\begin{somme}
    海  & \textit{sea}  & うみ (umi)  \\\hline
    海水  & sea + water = \textit{sea water}  & カイ{\middel}スイ (KAI{\middel}SUI)  \\\hline
    海図  & sea + diagram = \textit{nautical chart}  & カイ{\middel}ズ (KAI{\middel}ZU)  \\\hline
    海王星  & sea + king + star = \textit{Neptune (planet)}  & カイ{\middel}オウ{\middel}セイ
    (KAI{\middel}\={O}{\middel}SEI)  \\\hline
    北海道  & north + sea + road = \textit{Hokkaido}  & ホッ{\middel}カイ{\middel}ドウ
    (HOK{\middel}KAI{\middel}D\={O})  \\\hline
\end{somme}

\begin{sentence}
    \underline{\ruby{夏休}{なつやす}み}に\underline{\ruby{北海道}{ホッカイドウ}}\ruby{へ}{え}\underline{\ruby{行}{い}きました}。\\\hline
    natsu yasu{\middel}mi ni HOK{\middel}KAI{\middel}D\={O} e i{\middel}kimashita.\\\hline
    (summer vacation) (Hokkaido) (went)\\\hline
    =\textit{(We) went to Hokkaido for summer vacation.}\\\hline
\end{sentence}
\end{multicols}
\vspace*{-\baselineskip}
\vhrulefill{5pt}
%%--------------------------------------------------------------------
%14-------------------------------------------------------------------
\begin{multicols}{2}
\begin{glyph}{School (校)}{school}\label{glyph:school}
School.\\\hline
A secondary meaning related to ``proof-reading'' shows up in only a few
    compounds.
\end{glyph}

\begin{comps}
\comprtwocone & 木 + 交 \\\hline
    \comprthreecone & Tree (\ref{glyph:tree}) + Mix (\ref{glyph:mix}) \\\hline
\end{comps}

\begin{remglyph}
    A \hooglig{school} is an environment where growing \remcom{trees}
    \remcom{mix}.\\\hline
\end{remglyph}

\begin{prons}
    \pronsrtwocone & \textbf{コウ (K\={O})} & --- \\\hline
\pronsrthreecone & --- & --- \\\hline
\end{prons}

\begin{remprons}
    The \hooglig{school} must have a noble \comon{cause}.\\\hline
\end{remprons}

%{\middel}
% &  ()\\\hline
\begin{somme}
    \rye{3}{Notice how ``学'' in the second example doubles up with 校, and how
    this remains constant in the other compounds in which this word occurs.}
    校外  & school + outside = \textit{off-campus}  & コウ{\middel}ガイ (K\={O}{\middel}GAI) \\\hline
    学校  & study + school = \textit{school}  & ガッ{\middel}コウ (GAK{\middel}K\={O})  \\\hline
    小学校  & small + study + school = \textit{elementary school}  &
    ショウ{\middel}ガッ{\middel}コウ (SH\={O}{\middel}GAK{\middel}K\={O})  \\\hline
    中学校  & middle + study + school = \textit{junior high school}  &
    チュウ{\middel}ガッ{\middel}コウ (CH\={U}{\middel}GAK{\middel}K\={O})  \\\hline
    高校  & tall + school = \textit{high school}  & コウ{\middel}コウ (K\={O}{\middel}K\={O})  \\\hline
    女学校  & woman + study + school = \textit{girls' school}  & ジョ{\middel}ガッ{\middel}コウ
    (JO{\middel}GAK{\middel}K\={O})  \\\hline
\end{somme}

\begin{sentence}
    \underline{あの}{\;}\underline{\ruby{女}{おんな}}の\underline{\ruby{子}{こ}}\ruby{は}{わ}\underline{\ruby{高校}{コウコウ}}{\;}\underline{\ruby{二年生}{ニネンセイ}}です。\\\hline
    ano onna no ko wa K\={O}{\middel}K\={O} NI{\middel}NEN{\middel}SEI desu.\\\hline
    (that) (woman) (child) (high school) (second year student)\\\hline
    =\textit{That girl is a second-year high school student.}\\\hline
\end{sentence}
\end{multicols}
\vspace*{-\baselineskip}
\vhrulefill{5pt}
\vspace*{\fill}
\pagebreak
%%--------------------------------------------------------------------
%15-------------------------------------------------------------------
\vspace*{\fill}
\begin{multicols}{2}
\begin{glyph}{Leg (足)}{leg}\label{glyph:leg}
Confusingly for English speakers, this character can signify both ``leg'' and
    ``foot''.\\\hline
    A secondary meaning relates to the sense of ``suffice''.
\end{glyph}

\begin{comps}
\comprtwocone & 口 + 止 \\\hline
    \comprthreecone & Mouth (\ref{glyph:mouth}) + Stop (\ref{glyph:stop}) \\\hline
\end{comps}

\begin{remglyph}
    The \hooglig{leg} \remcom{stopped} \remcom{mouthing} to me.\\\hline
\end{remglyph}

\begin{prons}
\pronsrtwocone & \textbf{ソク (SOKU)} & \textbf{あし (ashi)}\\\hline
    \pronsrthreecone & --- & た (tunnel) \\\hline
\end{prons}

\begin{remprons}
    The \hooglig{leg} was \comkun{ushered} in by the \ucomkun{tunnel}
    \comon{sock}.\\\hline
\end{remprons}

%{\middel}
% &  ()\\\hline
\begin{somme}
    \rye{3}{Note that ``取'' is voiced in the fourth example.}
    足  & \textit{leg; foot} & あし (ashi)  \\\hline
    足首  & leg + neck = \textit{ankle}  & あし{\middel}くび (ashi{\middel}kubi)  \\\hline
    足音& leg + sound = \textit{sound of footsteps}  & あし{\middel}おと (ashi{\middel}oto)  \\\hline
    足取り& leg + take = \textit{one's gait}  & あし　ど{\middel}り (ashi do{\middel}ri)  \\\hline
    土足& earth + leg = \textit{with shoes on}  & ド{\middel}ソク (DO{\middel}SOKU)  \\\hline
    足元注意 & leg + basis + pour + mind = \textit{``watch your step''}  &
    あし{\middel}もと{\middel}チュウ{\middel}イ (ashi{\middel}moto{\middel}CH\={U}{\middel}I)  \\\hline
\end{somme}

\begin{irrr}
裸足  & naked + leg = \textit{barefoot}  & はだし (hadashi) \\\hline
\end{irrr}

\begin{sentence}
    \underline{\ruby{足音}{あしおと}}\ruby{を}{お}\underline{\ruby{聞}{き}きました}か。\\\hline
ashi{\middel}oto o ki{\middel}kimashita ka.\\\hline
    (sound of footsteps) (heard)\\\hline
    =\textit{Did you hear footsteps?}\\\hline
\end{sentence}
\end{multicols}
\vspace*{-\baselineskip}
\vhrulefill{5pt}
%%--------------------------------------------------------------------
%16-------------------------------------------------------------------
\begin{multicols}{2}
\begin{glyph}{Thread (糸)}{thread}\label{glyph:thread}
Thread.  This character appears as a component in more than sixty other kanji.
\end{glyph}

\begin{comps}
\comprtwocone & 幺 + 小 \\\hline
    \comprthreecone & small/young/short thread + Small (\ref{glyph:small})  \\\hline
\end{comps}

\begin{remglyph}
    The \hooglig{thread} is \remcom{young} and \remcom{small}.\\\hline
\end{remglyph}

\begin{prons}
\pronsrtwocone & --- & \textbf{いと (ito)}\\\hline
    \pronsrthreecone & シ (SHI) & --- \\\hline
\end{prons}

\begin{remprons}
    \hooglig{Threaded} \ucomon{shields} against \comkun{e-tolling}.\\\hline
\end{remprons}

%{\middel}
% &  ()\\\hline
\begin{somme}
    \rye{3}{Note that ``口'' is voiced in the second example.}
    糸  & \textit{thread}  & いと (ito)  \\\hline
    糸口  & thread + mouth = \textit{beginning; clue}  & いとぐち (ito{\middel}guchi)  \\\hline
\end{somme}

\begin{sentence}
    \underline{\ruby{青}{あお}い}{\;}\underline{\ruby{糸}{いと}}が\underline{あります}か。\\\hline
ao{\middel}i ito ga arimasu ka.\\\hline
    (blue) (thread) (is)\\\hline
    =\textit{Is there any blue thread?}\\\hline
\end{sentence}
\end{multicols}
\vspace*{-\baselineskip}
\vhrulefill{5pt}
\vspace*{\fill}
\pagebreak
%%--------------------------------------------------------------------
%17-------------------------------------------------------------------
\vspace*{\fill}
\begin{multicols}{2}
\begin{glyph}{Same (同)}{same}\label{glyph:same}
Same.
\end{glyph}

\begin{comps}
\comprtwocone & 冂 + 一 + 口 \\\hline
    \comprthreecone & wilderness + One (\ref{glyph:one}) + Mouth
    (\ref{glyph:mouth}) \\\hline
\end{comps}

\begin{remglyph}
    The \hooglig{same} \remcom{one} \remcom{vampire} in the \remcom{wilderness}.\\\hline
\end{remglyph}

\begin{prons}
    \pronsrtwocone & \textbf{ドウ (D\={O})} & \textbf{おな (ona)}\\\hline
    \rye{3}{Remember that おな differs from おんな, a reading for ``女'' in Entry
    16.}
\pronsrthreecone & --- & --- \\\hline
\end{prons}

\begin{remprons}
    \hooglig{Same} \comkun{onanism} \comon{doorway}.\\\hline
\end{remprons}

%{\middel}
% &  ()\\\hline
\begin{somme}
    同じ  & \textit{same}  & おな{\middel}じ (ona{\middel}ji)  \\\hline
    同意  & same + mind = \textit{agreement}  & ドウ{\middel}イ (D\={O}{\middel}I)  \\\hline
    同一  & same + one = \textit{identical}  & ドウ{\middel}イツ (D\={O}{\middel}ITSU)  \\\hline
    同時  & same + time = \textit{simultaneous}  & ドウ{\middel}ジ (D\={O}{\middel}JI)  \\\hline
    同化  & same + change = \textit{assimilation}  & ドウ{\middel}カ (D\={O}{\middel}KA)  \\\hline
    同行者  & same + go + individual = \textit{traveling companion}  &
    ドウ{\middel}カ{\middel}シャ (D\={O}{\middel}K\={O}{\middel}SHA)  \\\hline
\end{somme}

\begin{sentence}
    \underline{\ruby{私}{わたし}}と\underline{\ruby{中村}{なかむら}さん}\ruby{は}{わ}\underline{\ruby{同時}{ドウジ}}に\underline{\ruby{来}{き}ました}。\\\hline
    watashi to Naka{\middel}mura-san wa D\={O}{\middel}JI ni ki{\middel}mashita.\\\hline
    (I) (Nakamura-san) (simultaneous) (came)\\\hline
    =\textit{Nakamura-san and I arrived at the same time.}\\\hline
\end{sentence}
\end{multicols}
\vspace*{-\baselineskip}
\vhrulefill{5pt}
%%--------------------------------------------------------------------
%18-------------------------------------------------------------------
\begin{multicols}{2}
\begin{glyph}{Noon (午)}{noon}\label{glyph:noon}
Noon.
\end{glyph}

\begin{comps}
\comprtwocone & ノ + 干 \\\hline
    \comprthreecone & NO/slash + Dry (\ref{glyph:dry}) \\\hline
\end{comps}

\begin{remglyph}
    \hooglig{Noon} \remcom{drying} is a \remcom{no-no}.\\\hline
\end{remglyph}

\begin{prons}
\pronsrtwocone & \textbf{ゴ (GO)} & --- \\\hline
\pronsrthreecone & --- & --- \\\hline
\end{prons}

\begin{remprons}
    \hooglig{Noon} \comon{gobbling}.\\\hline
\end{remprons}

%{\middel}
% &  ()\\\hline
\begin{somme}
    午前  & noon + before = \textit{A.M.; morning}  & ゴ{\middel}ゼン (GO{\middel}ZEN)  \\\hline
    午後  & noon + after = \textit{P.M.; afternoon}  & ゴ{\middel}ゴ (GO{\middel}GO)  \\\hline
    午前中  & noon + before + middle = \textit{all morning}  & ゴ{\middel}ゼン{\middel}チュウ
    (GO{\middel}ZEN{\middel}CH\={U})  \\\hline
\end{somme}

\begin{sentence}
    \underline{\ruby{午前}{ゴゼン}}{\;}\underline{\ruby{八時}{ハチジ}}に\underline{\ruby{高校}{コウコウ}}\ruby{へ}{え}\underline{\ruby{行}{い}きます}。\\\hline
    GO{\middel}ZEN HACHI{\middel}JI ni K\={O}{\middel}K\={O} e i{\middel}kimasu.\\\hline
    (A.M.) (eight o'clock) (high school) (go)\\\hline
    =\textit{I go to high school at 8 A.M.}\\\hline
\end{sentence}
\end{multicols}
\vspace*{-\baselineskip}
\vhrulefill{5pt}
\vspace*{\fill}
\pagebreak
%%--------------------------------------------------------------------
%19-------------------------------------------------------------------
\vspace*{\fill}
\begin{multicols}{2}
\begin{glyph}{Love (愛)}{love}\label{glyph:love}
Love.
\end{glyph}

\begin{comps}
\comprtwocone & 爪 + 冖 + 心 + 夊 \\\hline
    \comprthreecone & claw + cover/roof + Heart (\ref{glyph:heart}) + go \\\hline
\end{comps}

\begin{remglyph}
    \hooglig{Love} from the \remcom{heart} \remcom{goes} through the
    \remcom{clawed} \remcom{roof}.\\\hline
\end{remglyph}

\begin{prons}
\pronsrtwocone & \textbf{アイ (AI)} & --- \\\hline
\pronsrthreecone & --- & --- \\\hline
\end{prons}

\begin{remprons}
    \hooglig{Love} melts the \comon{ice} away.\\\hline
\end{remprons}

%{\middel}
% &  ()\\\hline
\begin{somme}
    愛好  & love + like = \textit{love; like}  & アイ{\middel}コウ (AI{\middel}K\={O})  \\\hline
    愛犬  & love + dog = \textit{pet dog}  & アイ{\middel}ケン (AI{\middel}KEN)  \\\hline
    愛犬家  & love + dog + house = \textit{dog lover}  & アイ{\middel}ケン{\middel}カ (AI{\middel}KEN{\middel}KA)  \\\hline
    愛鳥家  & love + bird + house = \textit{bird lover}  & アイ{\middel}チョウ{\middel}カ
    (AI{\middel}CH\={O}{\middel}KA)  \\\hline
    愛国心  & love + country + heart = \textit{patriotism}  & アイ{\middel}コク{\middel}シン
    (AI{\middel}KOKU{\middel}SHIN)  \\\hline
\end{somme}

\begin{irrr}
可愛い  & possible + love = \textit{cute}  & カ{\middel}ワイ{\middel}い (KA{\middel}WAI{\middel}i)  \\\hline
\end{irrr}

\begin{sentence}
    \underline{\ruby{本田}{ほんだ}さん}\ruby{は}{わ}\underline{\ruby{愛鳥家}{アイチョウカ}}{\;}\underline{として}{\;}\underline{\ruby{有名}{ユウメイ}}です。\\\hline
    Hon{\middel}da-san wa AI{\middel}CH\={O}{\middel}KA toshite Y\={U}{\middel}MEI desu.\\\hline
    (Honda-san) (bird lover) (as) (famouns)\\\hline
    =\textit{Honda-san is famous as a bird lover.}\\\hline
\end{sentence}
\end{multicols}
\vspace*{-\baselineskip}
\vhrulefill{5pt}
%%--------------------------------------------------------------------
%20-------------------------------------------------------------------
\begin{multicols}{2}
\begin{glyph}{Separate (離)}{separate}\label{glyph:separate}
Separate/Leave.
\end{glyph}

\begin{comps}
\comprtwocone & 亠 + 凶 + 冂 + 厶 + 隹 \\\hline
    \comprthreecone & head/above + Evil (\ref{glyph:evil}) + wilderness +
    self/private + small bird \\\hline
\end{comps}

\begin{remglyph}
    \hooglig{Separate} \remcom{evil} from the \remcom{heads} of \remcom{small
    birds} in the \remcom{private} \remcom{wilderness}.\\\hline
\end{remglyph}

\begin{prons}
\pronsrtwocone & \textbf{リ (RI)} & \textbf{はな (hana)}\\\hline
\pronsrthreecone & --- & --- \\\hline
\end{prons}

\begin{remprons}
    \hooglig{Separate} \comkun{hung-up} \comon{leaps}.\\\hline
\end{remprons}

%{\middel}
% &  ()\\\hline
\begin{somme}
    離れる (intr.)  & \textit{to separate; to leave}  & はな{\middel}れる (hana{\middel}reru)  \\\hline
    離す (tr.)  & \textit{to separate (something)}  & はな{\middel}す (hana{\middel}su)  \\\hline
    離れ離れ  & separate + separate = \textit{scattered; separated}  &
    はな{\middel}れ　ばな{\middel}れ (hana{\middel}re bana{\middel}re)  \\\hline
    切り離す (tr.)  & cut + separate = \textit{to cut off}  & き{\middel}り　はな{\middel}す (ki{\middel}ri
    hana{\middel}su)  \\\hline
    分離  & part + separate = \textit{separation}  & ブン{\middel}リ (BUN{\middel}RI)  \\\hline
    離村  & separate + village = \textit{rural exodus}  & リ{\middel}ソン (RI{\middel}SON)  \\\hline
\end{somme}

\begin{sentence}
    \underline{\ruby{山}{やま}}で\underline{その}{\;}\underline{\ruby{親子}{おやこ}}\ruby{は}{わ}\underline{\ruby{離}{はな}れ\ruby{離}{ばな}れ}に\underline{なってしまいました}。\\\hline
yama de sono oya{\middel}ko wa hana{\middel}re bana{\middel}re ni natte shimaimashita.\\\hline
    (mountain) (that) (parent and child) (separated) (ended up)\\\hline
    =\textit{That parent and child ended up getting separated on the
    mountain.}\\\hline
\end{sentence}
\end{multicols}
\vspace*{-\baselineskip}
\vhrulefill{5pt}
%%--------------------------------------------------------------------
%%--------------------------------------------------------------------
\vspace*{\fill}
\pagebreak
